\documentclass{standalone}

\usepackage[T2A]{fontenc}
\usepackage[utf8]{inputenc}
\usepackage[russian]{babel}
\usepackage{amsmath}
\usepackage{amsfonts}
\usepackage{amssymb}

\usepackage{pgfplots}
\pgfplotsset{compat=1.18}
\usepackage{tikz}

\pgfmathsetmacro{\C}{34.3e-9}      % 34.3 * 10^(-9) Ф
\pgfmathsetmacro{\L}{0.11}          % 0.11 Гн
\pgfmathsetmacro{\R}{0.63*1000}     % 630 Ом (обновлено)
\pgfmathsetmacro{\E}{0.21}           % 0.21 В

\begin{document}
	
	\begin{tikzpicture}
		\begin{axis}[
			width=16cm,
			height=10cm,
			xlabel={$w, \text{рад/с}$},
			ylabel={$U_r, \text{В}$},
			xmin=9000, xmax=32000,
			ymin=0.025, ymax=0.225,
			grid=both,
			minor x tick num=10,
			minor y tick num=15,
			minor grid style={gray!85,  thin},
			grid style={black!50, thin},
			major grid style={black!70, thin},
			major tick length=2pt,
			minor tick length=1pt,
			axis lines=left,
			axis line style={-stealth, thick, black},
			tick style={black, thick},
			xmajorgrids=true,
			ymajorgrids=true,
			xminorgrids=true,
			yminorgrids=true,
			xtick distance=1000,  % Деления через каждые 2000 единиц
			ytick distance=0.025,
			scale only axis
			]
			
			% Добавляем точки
			\addplot[only marks, mark=*, mark size=1pt, color=blue] coordinates {
				(9424, 0.0599)
				(10681, 0.076)
				(11938, 0.097)
				(12566, 0.11)
				(13823, 0.1409)
				(15079, 0.1748)
				(16336, 0.1957)
				(16399, 0.1961)
				(16430, 0.1962)
				(16461, 0.1963)
				(16493, 0.1965)
				(16524, 0.1965)
				(16556, 0.1966)
				(16587, 0.1966)
				(16619, 0.1966)
				(16650, 0.1967)
				(16718, 0.1967)
				(16744, 0.1965)
				(16776, 0.1965)
				(16838, 0.1962)
				(16901, 0.1961)
				(16964, 0.1956)
				(17027, 0.1952)
				(17090, 0.1947)
				(17153, 0.1942)
				(17215, 0.1936)
				(17247, 0.1932)
				(17278, 0.1929)
				(17310, 0.1925)
				(17341, 0.1923)
				(17467, 0.1906)
				(17592, 0.1887)
				(17718, 0.1869)
				(17969, 0.1824)
				(18346, 0.1751)
				(18723, 0.1673)
				(19792, 0.1455)
				(21991, 0.1104)
				(25761, 0.0764)
				(30787, 0.0542)
				
			};
			\addlegendentry{$U_r - \text{практика}$}
			
			% U_C (напряжение на конденсаторе)
			\addplot[
			domain=9000:32000,
			samples=300,
			smooth,
			thick,
			color=blue
			] 
			{ 0.21 * 630 / sqrt( 630^2 + (x * 0.18 - 1/(x * 34.3e-9) )^2 )};
			\addlegendentry{$U_r - \text{теория}$}
		\end{axis}
	\end{tikzpicture}
\end{document}